\documentclass[11pt]{article}

\usepackage[margin=1in]{geometry}
\usepackage{booktabs}
\usepackage{amsmath}
\usepackage{hyperref}

\title{Dual-Domain MMD Adaptation: Large-Source, Small-Target Runs vs.\ a No-MMD
Fine-Tune Control}
\author{}
\date{\today}

\begin{document}
\maketitle

\section{Setup}

Both directions of adaptation were run from a large (full-size) pretrained checkpoint,
fine-tuned against the small opposite-domain target set:

\begin{itemize}
    \item \textbf{real\_large\_source\_syn\_small\_target}: source = real\_large
    (pretrained, detached, forward-only reference), target = syn\_small (fine-tuned via
    detection loss, pulled toward source via MMD). 50 epochs.
    \item \textbf{syn\_large\_source\_real\_small\_target}: mirror direction, source =
    syn\_large, target = real\_small. 50 epochs.
\end{itemize}

For each direction, a no-MMD control was run as a plain continued fine-tune from the
same starting checkpoint, on the same small target set, with no MMD term and no source
batches at all:

\begin{itemize}
    \item \textbf{syn\_small\_from\_real\_large}: plain fine-tune of real\_large's
    checkpoint on syn\_small only. 100 epochs.
    \item \textbf{real\_small\_from\_syn\_large}: plain fine-tune of syn\_large's
    checkpoint on real\_small only. 100 epochs.
\end{itemize}

\paragraph{Exposure-count note.} Epoch length in the MMD runs is defined by the
\emph{source} domain (the large set), with the small target resampled randomly each
step; at batch size 8 this works out to roughly 200$\times$ oversampling of the
1600-image target set over 50 epochs. The no-MMD controls, training directly on the
small target set, get exactly 100$\times$ exposure over their 100 epochs. The MMD runs
therefore received \emph{more}, not less, total exposure to the target set than their
no-MMD counterparts -- so the comparison below is not biased in the no-MMD control's
favor.

\section{Results}

\begin{table}[h!]
    \centering
    \begin{tabular}{lrrrr}
        \toprule
        Model & person mAP50 & person mAP50-95 & vehicle mAP50 & vehicle mAP50-95 \\
        \midrule
        real\_large\_source\_syn\_small\_target (MMD) & 0.693 & 0.357 & 0.697 & 0.483 \\
        syn\_small\_from\_real\_large (no-MMD) & 0.686 & 0.354 & 0.693 & 0.484 \\
        \midrule
        syn\_large\_source\_real\_small\_target (MMD) & 0.111 & 0.0266 & 0.214 & 0.138 \\
        real\_small\_from\_syn\_large (no-MMD) & 0.129 & 0.0349 & 0.222 & 0.146 \\
        \bottomrule
    \end{tabular}
    \caption{Validated on the synthetic (syn) val set. Top pair: syn\_small is the
    target domain (the direct MMD-vs-no-MMD comparison of interest). Bottom pair:
    syn is the source (forgotten) domain for that direction.}
\end{table}

\begin{table}[h!]
    \centering
    \begin{tabular}{lrrrr}
        \toprule
        Model & person mAP50 & person mAP50-95 & vehicle mAP50 & vehicle mAP50-95 \\
        \midrule
        syn\_large\_source\_real\_small\_target (MMD) & 0.642 & 0.296 & 0.877 & 0.625 \\
        real\_small\_from\_syn\_large (no-MMD) & 0.644 & 0.303 & 0.877 & 0.626 \\
        \midrule
        real\_large\_source\_syn\_small\_target (MMD) & 0.028 & 0.00748 & 0.340 & 0.205 \\
        syn\_small\_from\_real\_large (no-MMD) & 0.0146 & 0.00418 & 0.361 & 0.219 \\
        \bottomrule
    \end{tabular}
    \caption{Validated on the real val set. Top pair: real\_small is the target domain
    (the direct comparison of interest). Bottom pair: real is the source (forgotten)
    domain for that direction.}
\end{table}

\section{Discussion}

\textbf{No visible MMD benefit on the target-domain class that matters most.} On the
metric of primary interest -- vehicle mAP on whichever domain is actually being
fine-tuned for -- the MMD and no-MMD runs are essentially indistinguishable in both
directions: 0.877/0.625 vs.\ 0.877/0.626 (real, direction 1) and 0.697/0.483 vs.\
0.693/0.484 (syn, direction 2). The tiny edge that exists flips sign between the two
directions (no-MMD marginally ahead in direction 1, MMD marginally ahead in direction
2), consistent with run-to-run noise rather than a systematic effect.

\textbf{This contrasts with the small-source/large-target regime.} In that earlier
setup (see \texttt{experiment\_summary\_v3.tex}), MMD showed a small but consistent
vehicle-mAP improvement over the target domain's own in-domain baseline, in both
directions. Here, with the roles reversed (large source, small target), that benefit
does not appear.

\textbf{A plausible explanation.} MMD's contribution may be tied to how much direct,
diverse target-domain data the model actually trains on. When the target set is large
and gets extensive direct exposure (small-source/large-target regime), MMD's
feature-space regularization has room to meaningfully shape the shared backbone. When
the target set is small and relies on heavy resampling with the same $\sim$1600 images
(large-source/small-target regime), plain fine-tuning on that same scarce set already
extracts what there is to extract, and the added MMD machinery does not provide a
further edge -- despite the MMD runs here having received \emph{more} total exposure to
the target set (via oversampling) than the no-MMD controls, not less.

\textbf{Caveats.} These are single runs on val sets of modest size (160--640 images);
the MMD runs used 50 epochs against the no-MMD controls' 100, and while the
oversampling arithmetic means MMD still saw more total target-image draws, epoch count
and optimizer-schedule details (learning-rate decay shape in particular) were not
otherwise matched. A repeat run, or matching total step count exactly rather than just
approximately, would be needed before treating "no MMD benefit in this regime" as a
settled result.

\end{document}
